\documentclass[11pt]{article}
\usepackage[T1]{fontenc}
\usepackage[utf8]{inputenc}
\usepackage{lmodern,microtype}
\usepackage[a4paper,margin=27mm]{geometry}
\usepackage{amsmath,amssymb,amsthm,mathtools,booktabs,array}
\usepackage{hyperref,xurl}
\hypersetup{colorlinks=true,linkcolor=blue,citecolor=blue,urlcolor=blue,
 pdftitle={All-real Jacobians and exact response fibers of finite tree compartment models},
 pdfauthor={Yongxian Zhang}}
\newtheorem{theorem}{Theorem}
\newtheorem{corollary}[theorem]{Corollary}
\newtheorem{lemma}[theorem]{Lemma}
\newtheorem{proposition}[theorem]{Proposition}
\theoremstyle{remark}
\newtheorem*{remark}{Remark}
\newcommand{\R}{\mathbb R}
\newcommand{\C}{\mathbb C}
\newcommand{\Res}{\operatorname{Res}}
\newcommand{\diag}{\operatorname{diag}}
\newcommand{\1}{\mathbf 1}
\newcommand{\coef}{\operatorname{coef}}
\DeclareMathOperator{\supp}{supp}
\setlength{\emergencystretch}{2em}
\title{All-real Jacobians and exact response fibers\newline
of finite tree compartment models}
\author{Yongxian Zhang\\
\small School of Computer Science and Engineering\\
\small South China University of Technology\\
\small \href{mailto:mxymmxym1@gmail.com}{mxymmxym1@gmail.com}\\
\small ORCID: \href{https://orcid.org/0009-0000-3864-3536}{0009-0000-3864-3536}}
\date{14 October 2026}
\begin{document}
\maketitle
\begin{abstract}
For a finite labelled bidirected tree with collocated input and state
observation at $o$ and one external leak at any vertex $\ell$, we factor the
actual characteristic/cofactor coefficient Jacobian at every real parameter
vector. The factors are one rate on each edge directed toward $\ell$, each
edge product to its observed-subtree-size-minus-one power, and the squared
resultants of the actual sibling principal polynomials. Signed triangular
witnesses prove every coordinate multiplicity is exact, including multiplicity
zero for the leak, and identify the entire real Jacobian zero set.
The root-leak specializations imply the catenary and tree-multiplicity
conjectures of Gross, Meshkat and Shiu. For every full positive Stieltjes
response, we also classify all real realizations: positivity and Jacobian
regularity are consequences of the response, and the exact count is
$(N-1)!/\prod_{v\ne o}|T_v|$ at each prescribed leak site. An unknown site
multiplies this actual matrix count by $N$. We supply ordinary proofs and a
map to a fixed, independently kernel-checked Lean package. Verification
certifies the specified statements; historical priority and external peer
review are separate questions.
\end{abstract}

\section{Context and precise contribution}
Identifiability of linear compartmental systems has a substantial classical
literature. Tree-type realization counts were already investigated by Cobelli,
Lepschy and Romanin Jacur \cite{CLR}; we make no claim to have invented tree
enumeration or Stieltjes reconstruction. Gross, Meshkat and Shiu
\cite{GMS} posed a path singular-locus monomial conjecture (Conjecture 5.5)
and a rooted-tree coordinate-multiplicity conjecture (Conjecture 5.9), the
latter attributed there to Hoch, Sweeney and Tung. Our factorization and
exact-power statements imply both in their original collocated-leak setting.
Ahmed et al.\ \cite{Ahmed}, Section 6, still discuss the earlier path
conjecture as a prospective application of their coefficient formulas.
These correspondences do not establish that the present proof is historically
first: targeted searches cannot rule out all later or unpublished work.

The paper's main algebraic result includes arbitrary single-leak location,
every real rate, the full sibling-resultant residual, and exact coordinate
powers. A second result gives an exhaustive real fiber for each full
distinct-positive-pole target, with automatic admissibility of every real
competitor. It distinguishes this non-generic target statement from a generic
complex identifiability degree. The matrix, coefficients, derivatives and
ODE observations are the actual compartment objects throughout.

\section{Actual model and all-real factorization}
Let $T=(V,E)$ be a finite labelled simple tree, $|V|=N\ge1$. Fix the
collocated input and state observer $o\in V$, and one leak vertex $\ell\in V$.
For each directed edge $u\to v$ let $k_{v u}\in\R$ be its transfer rate,
and let $\lambda\in\R$ be the external leak at $\ell$. Define
\begin{equation}\label{eq:model}
 A_{v u}=k_{v u}\quad(u\sim v),\qquad
 A_{u u}=-\sum_{v\sim u}k_{v u}-\lambda\,\1_{u=\ell},
\end{equation}
with all other off-diagonal entries zero. Thus
$\1^T A=-\lambda e_\ell^T$. The dynamics are
$\dot x=Ax+e_o u$, with state output $y=e_o^Tx$.

Put
\begin{equation}\label{eq:pair}
 P(s)=\det(sI-A),\qquad
 F_o(s)=\det((sI-A)[V\setminus\{o\}]),\qquad
 G_o(s)=\frac{F_o(s)}{P(s)}.
\end{equation}
The principal submatrix retains its original diagonal entries, including
rates exiting the retained set. The empty determinant is one. The two
polynomials are monic of degrees $N$ and $N-1$. No cancellation is performed
before defining the coefficient map.

Root the tree at $o$. For $v\ne o$, let $p(v)$ be its parent,
$T_v$ its descendant vertex set, $N_v=|T_v|$,
$a_v=k_{p(v),v}$ and $f_v=k_{v,p(v)}$. Fix a parent-before-child enumeration
and use parameter order $(\lambda,(a_v)_{v\ne o},(f_v)_{v\ne o})$.
The coefficient output is
\[
 C_{o,\ell}=(P_0,\ldots,P_{N-1},(F_o)_0,\ldots,(F_o)_{N-2}),
\]
in ascending powers; monic leading coefficients are omitted. This is a square
polynomial map in $2N-1$ real variables. Write
$J_{o,\ell}=\det DC_{o,\ell}$.
The Lean graph entry uses the same ascending array followed by a fixed
coordinate permutation; transporting between conventions changes only a
rate-independent sign. It never introduces a parameter-dependent factor.

For every child $c$ of any observed-root vertex $v$, define the actual
principal polynomial
\[
 P_c(s)=\det(sI-A[T_c]).
\]
Define the observed sibling residual and leak-oriented product by
\begin{equation}\label{eq:residual}
 S_{o,\ell}=\prod_{v\in V}\ \prod_{\{c,d\}\subseteq\operatorname{child}(v)}
       \Res(P_c,P_d)^2,\qquad
 L_\ell=\prod_{e\in E} k_{e\to\ell}.
\end{equation}
Here $k_{e\to\ell}$ is the rate on $e$ pointing toward $\ell$ along its
unique tree path. Empty products are one. Each $P_c$, including its diagonal,
depends on the original model with leak at $\ell$.

\begin{theorem}[All-real Jacobian and exact powers]\label{thm:jac}
There is a sign $\varepsilon_o\in\{1,-1\}$, depending only on the chosen
observation-root coordinate convention, such that simultaneously for every
leak site $\ell$,
\begin{equation}\label{eq:main}
 J_{o,\ell}=\varepsilon_o L_\ell
       \prod_{v\ne o}(a_vf_v)^{N_v-1}S_{o,\ell}.
\end{equation}
This is an identity of multivariate real polynomials and of the actual
Fr\'{e}chet derivative determinant at every real vector. The exact power of
either directed coordinate on the edge at $v$ is
\begin{equation}\label{eq:exponent}
 N_v-1+\1_{\text{that direction points toward }\ell}.
\end{equation}
The leak coordinate has exact power zero. Hence $J_{o,\ell}=0$ precisely
when a rate with positive displayed exponent vanishes, or $S_{o,\ell}=0$.
Zero leak, zero edge rates, repeated spectra and signed parameters are included.
\end{theorem}

Exact power means that the stated power of the coordinate variable divides
the formal Jacobian polynomial and the next power does not. The residual
is not asserted irreducible. The zero set is a Jacobian-rank statement;
it need not characterize every notion of local uniqueness at singular points.

\begin{corollary}[Rooted trees and endpoint paths]\label{cor:path}
For $\ell=o$, the inward and outward coordinate powers are respectively
$N_v$ and $N_v-1$. On the path $1,\ldots,N$ with $o=\ell=1$,
\[
 J=\varepsilon_o\prod_{i=1}^{N-1}
       k_{i,i+1}^{N-i}k_{i+1,i}^{N-i-1}.
\]
These give the path monomial and rooted-tree multiplicities of
\cite[Conjectures 5.5 and 5.9]{GMS}, in the stated original coordinate
convention up to sign.
\end{corollary}

For the path there are no sibling pairs, so $S=1$. The rooted-tree powers
satisfy the subtree recursion $N_v=1+\sum_{c\in\operatorname{child}(v)}N_c$,
with the corresponding inward/outward leaf values $1,0$. This is the required
multiplicity implication, rather than mere divisibility by a candidate monomial.

\section{Proof of the polynomial theorem}
\subsection{Coefficient assembly and the root-leak case}
The following elementary determinant calculation is the algebraic source
of the squared sibling resultants.

\begin{lemma}[Polynomial gluing]\label{lem:glue}
Let $p_i$ be monic of degree $n_i\ge1$, let $q_i$ have degree less than $n_i$,
and let $d$ be a scalar. With $K=\sum_i n_i$, put
\[
 F=\prod_i p_i,\qquad
 P=(s+d)F-\sum_i q_i\prod_{j\ne i}p_j.
\]
The map from the non-leading coefficients of the $p_i$, all coefficients of
the $q_i$, and $d$, to the non-leading coefficients of $F,P$ has determinant
\[
 \sigma\prod_{i<j}\Res(p_i,p_j)^2,
\]
where $\sigma=\pm1$ depends only on the fixed coefficient ordering. This
identity holds without coprimality or simple-root hypotheses.
\end{lemma}
\begin{proof}
Write $\mathcal L_p(h_i)=\sum_i h_i\prod_{j\ne i}p_j$, with
$\deg h_i<n_i$. This is a square linear map to polynomials of degree below $K$.
The derivative of $(p_i)\mapsto F$ is exactly $\mathcal L_p$.
In coefficient bases its determinant is a fixed sign times
$\prod_{i<j}\Res(p_i,p_j)$.

Here is a direct verification of the latter assertion. Initially take all
roots to be simple and mutually disjoint in $\C$. Evaluate at the roots of
each $p_i$, in fixed block order. Only $h_i$ contributes at its own block,
with diagonal multipliers $\prod_{j\ne i}p_j(\alpha)$.
The quotient of the full Vandermonde determinant by the within-block
Vandermondes is the cross-root difference product. Cancelling one copy
of that product in the evaluation identity leaves one resultant for each
pair, with a fixed ordering sign. Both sides are polynomials in the
coefficients, so equality on this nonempty open set proves the identity
for every coefficient vector.

For fixed $(p_i)$, variations of $(q_i,d)$ give
$\delta P=\delta d\,F-\mathcal L_p(\delta q_i)$. Since $F$ is monic,
its $s^K$ coefficient supplies the $\delta d$ pivot; the remaining determinant
is the same $\det\mathcal L_p$ up to a fixed sign. The entire derivative is
block triangular because $F$ does not depend on $q_i,d$. Its determinant
therefore contains two copies of the resultant product. This also proves
the empty-family case $F=1,P=s+d$.
\end{proof}

Suppose first that the leak is at $o$. Its child branches have sizes $n_i$,
actual characteristic polynomials $P_i$ and root cofactors $F_i$.
The principal branch is a genuine child model whose leak is the inward
connector $a_i$. Original diagonals are essential here. Schur expansion,
as a polynomial identity, gives
\begin{equation}\label{eq:recursion}
 F_o=\prod_i P_i,\qquad
 P=(s+d)\prod_i P_i-\sum_i a_if_iF_i\prod_{j\ne i}P_j,
 \quad d=\lambda+\sum_i f_i.
\end{equation}
It can first be checked where the branch blocks are invertible and then
extended polynomially; it asserts an identity of actual determinants.

Keep the child rate vectors, including their leak $a_i$, and change the
remaining variables from $(\lambda,(f_i))$ to $(d,(a_if_i))$.
The determinant is $\prod_i a_i$, by its block triangular derivative,
at all rate values. Scaling the monic $F_i$ by $b_i=a_if_i$ changes
$(b_i,\coef F_i)$ into $\coef(b_iF_i)$ with determinant
$b_i^{n_i-1}$ up to a fixed sign, even at $b_i=0$.
Applying Lemma~\ref{lem:glue} and each actual child Jacobian gives
\begin{equation}\label{eq:Jrec}
 J_T=\sigma_T\prod_i a_i(a_if_i)^{n_i-1}J_{T_i}
             \prod_{i<j}\Res(P_i,P_j)^2.
\end{equation}
All coordinate changes here are explicitly described; no inverse rate is used.
Induction, starting with $P=s+\lambda,F=1,J=1$, flattens
\eqref{eq:Jrec} into \eqref{eq:main} for $\ell=o$.
The sign depends only on the rooted enumeration and the fixed permutations.

\subsection{Leak relocation and removal of denominators}
On the positive orthant, move a leak $\lambda$ at $v$ across an adjacent edge
to $u$. Put $a=k_{u v}$ and $f=k_{v u}$. Replace only these three coordinates by
\begin{equation}\label{eq:slide}
 a'=a+\lambda,\qquad f'=\frac{af}{a+\lambda},\qquad
 \lambda'=\frac{f\lambda}{a+\lambda},
\end{equation}
with the new leak at $u$. The diagonal entries stay the same, since
$f'+\lambda'=f$ and $a'=a+\lambda$.
Cutting that edge gives two components. A diagonal scaling constant on each
component, with ratio $D_v/D_u=a/(a+\lambda)$, gives $A'=DAD^{-1}$.
It preserves every principal characteristic polynomial and every edge product.
In particular it preserves both $P,F_o$ and the entire core
$W=\prod_{v\ne o}(a_vf_v)^{N_v-1}S_{o,\ell}$.

A direct derivative of \eqref{eq:slide}, in order $(\lambda,a,f)$ and
$(\lambda',a',f')$, has determinant $f/(a+\lambda)=\lambda'/\lambda$.
Compose the slides on the unique path from $\ell$ to $o$. If $R$ is the
resulting positive rate transformation and $\lambda_o$ its new leak,
\[
 C_{o,\ell}=C_{o,o}\circ R,\qquad
 \det DR=\lambda_o/\lambda.
\]
The root-leak factorization and the chain rule imply
\begin{equation}\label{eq:transport}
 J_{o,\ell}=\varepsilon_o\frac{\lambda_o}{\lambda}F_o^{\rm new}(0)W.
\end{equation}

For any single-leak tree model, conservation and expansion in the leak give
\begin{equation}\label{eq:cofactorzero}
 P(0)=\lambda F_\ell(0),\qquad F_\ell(0)=L_\ell.
\end{equation}
For completeness, the leak-free matrix has zero determinant because its
columns sum to zero, and the leak enters a single diagonal linearly.
For the second identity root at $\ell$: eliminating a nonroot leaf in the
positive model contributes its inward rate to the principal determinant,
and its Schur correction cancels its outward contribution to the parent
diagonal. Iterating gives exactly $L_\ell$. Since these are polynomial
identities, they hold at every real specialization as well.

The characteristic polynomial is invariant under $R$, so
$\lambda_oF_o^{\rm new}(0)=\lambda F_\ell(0)$.
Substitute into \eqref{eq:transport}. The result is
$J_{o,\ell}=\varepsilon_oF_\ell(0)W$, with no transformed rate or denominator.
Both sides are polynomial functions of the original coordinates. Equality
on the product of positive half-lines proves polynomial equality, by applying
the one-variable infinite-zero principle successively to each coordinate.
Thus \eqref{eq:main} holds on all of $\R^{2N-1}$, with the same sign for
every $\ell$. We have not divided at a zero-boundary point or asserted that
the positive orthant is dense in the entire parameter space.

\subsection{Exact coordinate powers}
Fix a rate coordinate, or the leak coordinate. Root the tree instead at
$\ell$, set all rates pointing away from $\ell$ to zero, and assign numbers
$b_v$ to the inward rates at nonroot vertices and $b_\ell$ to the leak.
Choose all $b_v$ distinct. Then $A$ is triangular in leak-root order, with
diagonal $-b_v$. Every principal characteristic polynomial is the product
of the diagonal factors for its original vertices, even when that principal
set is defined in observation-root order. Sibling sets are disjoint;
therefore all their resultants are nonzero, so $S_{o,\ell}\ne0$.

This witness can put the selected coordinate on its zero hyperplane.
An outward-from-leak rate is already zero. For an inward coordinate at $v$,
set $b_v=0$ and choose every other $b_w$ distinct and nonzero.
For the leak coordinate set $b_\ell=0$ instead. Signed choices are allowed
because the claim is algebraic. Thus no coordinate variable divides the
residual polynomial. Each coordinate variable is prime in the multivariate
polynomial ring, so its exact power in \eqref{eq:main} is the displayed
monomial exponent. The leak has exponent zero. Since the sign is nonzero,
the stated real zero set follows directly from the product. This finishes
the proof of Theorem~\ref{thm:jac}.

\section{Full positive targets and exhaustive real fibers}
Fix
\begin{equation}\label{eq:target}
 G(s)=\sum_{i=1}^N\frac{w_i}{s+r_i},\qquad
 r_i>0\text{ distinct},\quad w_i>0,\quad\sum_iw_i=1.
\end{equation}
Equality of responses means equality on the common real nonpole domain,
equivalently rational-function equality. A real competitor has the fixed
labelled tree $T$ and one prescribed leak site, with no sign, stability,
minimality or Jacobian condition imposed beforehand.

\begin{theorem}[Complete real full-response classification]\label{thm:fiber}
For each prescribed $\ell$, every real realization of \eqref{eq:target}
has strictly positive directed rates and leak and a nonzero actual
coefficient Jacobian. Its entire fiber is in bijection with recursive labelled
pole allocations and has cardinality
\begin{equation}\label{eq:count}
 D(T,o)=\frac{(N-1)!}{\prod_{v\ne o}N_v}.
\end{equation}
If the leak site is unknown, the fiber contains exactly $N D(T,o)$ distinct
physical matrices. Labels are retained; tree automorphisms are not quotiented.
The singleton case has one realization. These conclusions hold for every
target \eqref{eq:target}, not only a generic subset of that class.
\end{theorem}

\subsection{Reciprocal data and positive reconstruction}
Order $r_1<\cdots<r_N$. On each interval between consecutive negative poles,
$G'(s)=-\sum_i w_i/(s+r_i)^2<0$; its one-sided limits have opposite signs.
Thus its numerator has exactly $N-1$ simple negative zeros $-\beta_j$.
Partial fractions give the unique identity
\begin{equation}\label{eq:reciprocal}
 \frac1{G(s)}=s+d-\sum_{j=1}^{N-1}\frac{\gamma_j}{s+\beta_j},
 \qquad \beta_j>0,\quad\gamma_j=-\frac1{G'(-\beta_j)}>0.
\end{equation}
The monic degrees determine the linear polynomial part. In particular
\[
 \lambda=d-\sum_j\gamma_j/\beta_j=1/G(0)>0.
\]
For root-child sizes $n_c$, allocate the $N-1$ reciprocal poles into labelled
disjoint sets $I_c$ of those sizes. Put
\[
 b_c=\sum_{j\in I_c}\gamma_j,\qquad
 G_c(s)=\sum_{j\in I_c}\frac{\gamma_j/b_c}{s+\beta_j}.
\]
Recursively realize $G_c$ on its labelled child tree. Its root leak is
$a_c=1/G_c(0)>0$. Choose $f_c=b_c/a_c=\sum_{j\in I_c}\gamma_j/\beta_j$.
The parent leak is $1/G(0)$, and its full root diagonal is
$\lambda+\sum_c f_c=d$. Formula~\eqref{eq:recursion} recovers exactly the
target characteristic polynomial and cofactor. This constructs a positive
realization for every recursive allocation.

\subsection{Reverse extraction and the entire root-leak fiber}
Write the target monic denominator and numerator as $P,F$.
They are coprime: at the pole $-r_i$, the numerator is
$w_i\prod_{j\ne i}(r_j-r_i)\ne0$.
For any $N$-vertex real competitor with equal response, cross multiplication
gives $F_xP=FP_x$. Coprimality and the equal monic denominator degrees force
$P_x=P$ and $F_x=F$. Coefficient equality is therefore recovered from the
original response, not additional observed information.

When the leak is at $o$, its actual child denominators are monic and multiply
to $F=\prod_j(s+\beta_j)$. They must partition these distinct linear factors
in exactly one labelled root allocation. In \eqref{eq:recursion}, compare
their branch numerators with those of the positive reference for that allocation.
Reducing the numerator difference modulo each $P_c$ isolates its term;
the product of the other, coprime denominators is invertible modulo $P_c$.
The term has degree below $n_c$, so it must be zero. Thus
$a_cf_cF_c$ equals the reference branch numerator. Its leading coefficient
gives $a_cf_c=b_c>0$ and its monic cofactor is the reference cofactor.
Each actual child consequently has the full target $G_c$. Induction forces
its real rates and its leak $a_c$ positive; then $f_c=b_c/a_c>0$, and
the parent leak is $1/G(0)>0$. This is a proof about every real competitor.

Forward assembly and reverse extraction are inverse on actual rate vectors.
There are $(N-1)!/\prod_c n_c!$ labelled root allocations; permutations within
a block do not make new allocations. Hence
\[
 D(T,o)=\frac{(N-1)!}{\prod_c n_c!}\prod_c D(T_c,c).
\]
Induction gives \eqref{eq:count}, with finiteness established by the explicit
bijection. Different sibling branches receive disjoint poles at each depth,
so every sibling resultant is nonzero. Theorem~\ref{thm:jac} now gives
actual Jacobian regularity for all these realizations.

\subsection{Automatic positivity at an arbitrary leak site}
One must not move the leak of a signed competitor using positive-only
formulas before proving admissibility. We instead use the following
diagonal-gauge construction for general real tree-supported matrices.

Let $M$ have this tree support and the full pair $P,F$ at $o$; no conservation
or sign condition is imposed. The same child factor partition and
partial-fraction uniqueness recover its root diagonal, child cofactors,
and every boundary edge product $M_{oc}M_{co}=b_c>0$.
Inductively each actual child matrix is diagonally similar to a positive
child realization, with a nonzero real diagonal gauge normalized at its root.
Attach the positive reference as above. On child component $c$, multiply
its normalized gauge by $M_{co}/f_c$, which is nonzero because the edge
product is positive. The outward boundary entry then matches $M_{co}$;
the matching positive edge product forces the inward entry to match as well.
The child blocks and root diagonal already match. Thus
\begin{equation}\label{eq:gauge}
 M=DBD^{-1},\qquad D_{oo}=1,\quad D_{vv}\ne0,
\end{equation}
where $B$ is a positive conservative root-leak realization. The singleton
is immediate. This proves the construction for all tree sizes and permits
signed gauges.

For the actual competitor $A$ with leak at $\ell$, relocate the positive
reference $B$ from $o$ to $\ell$ using \eqref{eq:slide}; call the resulting
positive matrix $B_\ell$. It has the same full pair and is invertible
because $P(0)=\prod_i r_i>0$. Equation~\eqref{eq:gauge} and the positive
relocation give $A=H B_\ell H^{-1}$ with every diagonal $h_v$ nonzero.
Conservation determines this gauge. Indeed, if the two leaks are
$\lambda_A,\lambda_B$ at the same $\ell$, their column sums imply
\[
 h^T B_\ell=-\lambda_A h_\ell e_\ell^T,\qquad
 \1^TB_\ell=-\lambda_B e_\ell^T.
\]
Invertibility and $\lambda_B>0$ show
$h^T=(\lambda_Ah_\ell/\lambda_B)\1^T$.
Its $\ell$ coordinate, with $h_\ell\ne0$, gives
$\lambda_A=\lambda_B$ and $h$ constant. Therefore $A=B_\ell$.
Every arbitrary-leak real competitor is positive, with no division by an
unproved nonzero competitor rate or leak.

For positive models, relocation along a fixed path is invertible: the reverse
single-edge slide recovers all three coordinates in \eqref{eq:slide}.
It gives a bijection between the root-leak fiber and each prescribed-site
fiber. Principal resultants, edge products and positive rates preserve
regularity. Finally, distinct sites yield distinct actual matrices: the
column-sum vector $-\lambda e_\ell^T$, with $\lambda>0$, recovers the site
and leak, while off-diagonal entries recover all rates. Summing the $N$
disjoint fixed-site fibers proves the unknown-site count.

\section{Actual continuous-input observations}
For the target \eqref{eq:target}, write
\[
 g_G(t)=\sum_{i=1}^N w_i e^{-r_i t},\qquad
 y_G(t)=\int_0^t g_G(t-\tau)u(\tau)\,d\tau\quad(t\ge0).
\]
The target zero-initial output below means precisely $y_G$.
\begin{proposition}[History interpretation]\label{prop:history}
For the tree and target of Theorem~\ref{thm:fiber}, equality with the target
zero-initial output for every continuous real input is equivalent to
equality of the collocated transfer response. The same exhaustive real
fiber and site counts therefore apply to these actual ODE histories.
Positive leak relocation preserves every zero-initial collocated history
at any fixed observer. This preservation itself does not need a full-mode
target.
\end{proposition}
\begin{proof}
For any real finite matrix $A$ and continuous input, the actual solution is
\[
 x(t)=e^{tA}x(0)+\int_0^t e^{(t-\tau)A}e_o u(\tau)\,d\tau.
\]
This exists globally and is the unique classical solution. The observed
zero-initial kernel is $g_A(t)=e_o^Te^{tA}e_o$.
Equal transfer functions imply equal kernel moments at zero by expanding
the resolvent at infinity. Since the matrix exponential is entire, all
moments determine the same kernel and hence every convolution output.
Conversely equal outputs for all continuous inputs include the constant
input one; differentiating its integral output gives equal kernels for
$t\ge0$. Their Laplace transforms agree for sufficiently large real $s$,
giving equal rational responses. This reasoning requires no competitor
stability or positivity and has actual solutions, so the quantified history
predicate is not vacuous. Theorem~\ref{thm:fiber} applies.

For relocation, normalize its positive diagonal similarity by $D_{oo}=1$.
Then $D e_o=e_o$ and $e_o^TD=e_o^T$. It transports the forced equation,
zero initial state and observed coordinate exactly, preserving every history.
\end{proof}

State observation at $o$ is distinct from measuring discharge
$\lambda x_\ell$ at another leak vertex. The separately formalized
collocated-discharge and signed-control results are not additional main
claims here. In particular this paper does not assert arbitrary endpoint
steering using nonnegative rainfall alone.

\section{Proof correspondence, verification and limits}
The accompanying \texttt{THEOREM\_MAP.md} was prepared before this manuscript.
It maps every main statement and necessary bridge to existing source
declarations; no new theorem root is introduced by publication preparation.
Table~\ref{tab:map} gives compact navigation. Actual Lean types, including
their hypotheses, are in the linked source and entrypoint index.

\begin{table}[ht]
\centering\small
\begin{tabular}{@{}p{.28\linewidth}p{.66\linewidth}@{}}
\toprule
Paper scope & Existing module families\\
\midrule
Theorem 1: all-real identity & LocatedAllRateJacobian; LocatedPolynomialWeight;
LocatedExactMultiplicities; LocatedBoundarySingularSet\\[3pt]
Exact powers and boundary & GraphResidualWitness; CoordinateMultiplicity\\[3pt]
Corollary 2: original conjectures & CatenaryFormalPolynomial;
FormalTreePolynomial; LocatedExactMultiplicities\\[3pt]
Lemma 3 and root induction & FamilyGluing; FamilyResultants; RateChart;
TreeJacobianRecurrence; AllTreeJacobian\\[3pt]
Theorem 4: real fibers & RealFiberClassification; SupportedPositiveGauge;
ConservationGaugeRigidity; RealLocatedClassification\\[3pt]
Proposition 5: actual histories & RealLocatedODEClassification;
MatrixTimeIdentification; LocatedLeakDynamics\\
\bottomrule
\end{tabular}
\caption{Navigation into the fixed existing Lean package, not a new audit.}
\label{tab:map}
\end{table}

The mathematical and public-scope source is fixed at
\path{c2a2ce36ea6c5e000225574835f5a3ceddc748b4}.
Lean is pinned to 4.34.1 (commit
\path{5045d0056413266e57c625dcd7c365b10e377c52}); Mathlib is
\path{d13f23b723b8a846827a245b89c10fc7d3f11612}.
The original independent GitHub Actions run
\href{https://github.com/mxym/math/actions/runs/38383557356}{38383557356}
reports 220 owned modules, 1,677 exported roots and 53,370 used declarations,
with trust-level-zero empty-kernel replay. Its allowed axioms are
\texttt{Classical.choice}, \texttt{Quot.sound} and \texttt{propext};
91 deliberately false source mutations were rejected.
It starts with zero owned objects and builds the owned source afresh;
this is not a claim that all of Mathlib was recompiled.

On 14 October 2026, publication preparation independently compared all
235 proof, reproduction and scope inputs against the current tree, the
original source commit and the verified-source archive; all matched.
The original ZIP digest was also matched to the GitHub API artifact
metadata, and the corresponding job was confirmed successful.
This is a source/archive/API readback of an existing independent run,
not a new build or a newly expanded kernel audit.
The preserved package provides its exact logs, lock manifests, negative
controls and reproduction script. Our manuscript review is an internal
assistant review, not external human peer review.

The full-target assumptions exclude repeated poles, fewer-mode data and
complex-parameter fibers. The algebraic Jacobian theorem permits signed
and zero rates but retains the finite single-leak tree and collocated
input/state-observation model. General cyclic graphs, multiple leaks,
off-diagonal input/output, noise, nonlinear catchments and predictive
performance are outside this manuscript. The results do not determine
unknown tree topology, nor claim a new topology-count corollary.

\clearpage
\begin{samepage}
\paragraph{Data and proof availability.}
The fixed proof package and receipt are available in
\href{https://github.com/mxym/math}{the author's repository}, under
\path{formalizations/hydrology-tree-core} and
\path{verification/hydrology-tree-core-c2a2ce36}.
This paper's source, theorem map, review and preparation checks accompany
it under \path{preprints/finite-tree-jacobian-2026-10}.
The paper restates and organizes the existing finite-tree result, rather
than attributing that result to a new research round.
\end{samepage}

\paragraph{Declarations.}
The author received no external funding. OpenAI AI assistants assisted
research, formalization, writing and internal review. Original material:
copyright 2026 Yongxian Zhang, all rights reserved. Existing repository
and third-party notices remain in force. This prepared manuscript does
not itself assert an aiXiv submission, Release or DOI.

\begin{thebibliography}{9}
\bibitem{CLR}
C. Cobelli, A. Lepschy and G. Romanin Jacur,
Identifiability results on some constrained compartmental systems,
\emph{Mathematical Biosciences} \textbf{47} (1979), 173--195.
\href{https://doi.org/10.1016/0025-5564(79)90036-1}{doi:10.1016/0025-5564(79)90036-1}.
The author-repository abstract was checked; a complete copy of this article
was not obtained for the present comparison.

\bibitem{GMS}
E. Gross, N. Meshkat and A. Shiu,
Identifiability of linear compartmental models: the singular locus,
\emph{Advances in Applied Mathematics} \textbf{133} (2022), 102268.
\href{https://doi.org/10.1016/j.aam.2021.102268}{doi:10.1016/j.aam.2021.102268};
author version \href{https://arxiv.org/abs/1709.10013v3}{arXiv:1709.10013v3}.

\bibitem{Ahmed}
S. Ahmed, N. Crepeau, P. R. Dessauer Jr., A. Edozie, O. Garcia-Lopez,
T. Grimsley, J. Lopez Garcia, V. Neri and A. Shiu,
Identifiability of directed-cycle and catenary linear compartmental models,
\emph{SIAM Journal on Applied Dynamical Systems} \textbf{25} (2026), 304--350.
\href{https://doi.org/10.1137/25M1728636}{doi:10.1137/25M1728636};
author version \href{https://arxiv.org/html/2412.05283v3}{arXiv:2412.05283v3}.
\end{thebibliography}
\end{document}
